% =====================================================================
%  Chapter 1: The path integral on the lattice
%  Gattringer & Lang, "Quantum Chromodynamics on the Lattice" (LNP 788)
%  Reconstructed from extract/ch01.txt; equations auto-numbered.
% =====================================================================

\chapter{The path integral on the lattice}

The basic tool for quantizing fields on the lattice is the Euclidean path
integral. Our first chapter is dedicated to the introduction of the path
integral formalism and to its interpretation. In order to develop the idea
without getting lost in technicalities we introduce the path integral for the
simplest case, a scalar field theory. We derive and discuss the two key
equations of lattice field theory. The first key equation is
\begin{equation}
  \begin{split}
    \ev{O_2(t)\,O_1(0)}
    &= \lim_{T\to\infty}\frac{1}{Z_T}\tr\!\big[e^{-(T-t)H}O_2\,e^{-tH}O_1\big]\\
    &= \sum_n \bbraket{0}{O_2}{n}\,\bbraket{n}{O_1}{0}\,e^{-tE_n},
  \end{split}
\end{equation}
where $Z_T$ is a normalization factor given by $Z_T = \tr[e^{-TH}]$. The
left-hand side of (1.1) is the Euclidean correlation function of two operators
$O_1$ and $O_2$, $H$ is the Hamiltonian of the system. On the right-hand side
the Euclidean correlator is expressed as a sum over eigenstates of the
Hamiltonian operator labeled by $n$. The terms in the sum contain matrix
elements of the operators $O_i$ taken between the vacuum $\ket{0}$ and the
physical states $\ket{n}$. These matrix elements are weighted with exponentials
containing the energy eigenvalues $E_n$ of the system. The right-hand side of
(1.1) can thus be used to extract matrix elements of operators as well as the
energy spectrum of the theory.

In the second key equation
\begin{equation}
  \frac{1}{Z_T}\tr\!\big[e^{-(T-t)H}O_2\,e^{-tH}O_1\big]
  = \frac{1}{Z_T}\int D[\Phi]\,e^{-S_E[\Phi]}\,O_2[\Phi(\cdot,t)]\,O_1[\Phi(\cdot,0)],
\end{equation}
the Euclidean correlator on the left-hand side is expressed as a path integral,
which is an integral over all possible configurations of the field $\Phi$. This
is a crucial point. The left-hand side is formulated in the operator language
of quantum field theory. The right-hand side knows nothing about field
operators. In the integrand the two operators $O_i$ are translated to
functionals $O_i$ of the fields and then weighted with the Boltzmann factor
containing the classical Euclidean action $S_E[\Phi]$. The right-hand side of
(1.2) can be evaluated numerically on the lattice.

\section{Hilbert space and propagation in Euclidean time}

This first section is dedicated to a detailed discussion of the Euclidean
correlators (1.1). Before we actually introduce the Euclidean correlators we
prepare the ground with a brief summary of the definition and properties of
Hilbert spaces (see, e.g., [1--4] for introductory texts).

\subsection{Hilbert spaces}

A Hilbert space $\mathcal{H}$ is an infinite dimensional vector space. Its
elements can be added and multiplied with scalars and $\mathcal{H}$ is closed
under these operations, i.e., for vectors $\ket{u}, \ket{v} \in
\mathcal{H}$ and complex numbers $\alpha$, $\beta$ we find
\begin{equation}
  \alpha\ket{u} + \beta\ket{v} \in \mathcal{H}.
\end{equation}
We will often refer to vectors in Hilbert space as states. In addition, a
Hilbert space is equipped with a scalar product, i.e., a sesquilinear
functional $\braket{u}{v}$ which maps a vector $\ket{v}$ and a dual vector
$\bra{u}$ into the complex numbers. The scalar product obeys the properties
(the $*$ denotes complex conjugation)
\begin{align}
  \braket{u}{v} &= \braket{v}{u}^*, \\
  \braket{w}{\alpha u + \beta v} &= \alpha\braket{w}{u} + \beta\braket{w}{v}.
\end{align}
The Hilbert spaces we are interested in here have a complete, countable basis,
i.e., a set of linearly independent vectors $\ket{e_n}$ such that every vector
$\ket{u} \in \mathcal{H}$ can be written as a linear combination of the basis
elements,
\begin{equation}
  \ket{u} = \sum_n \alpha_n \ket{e_n},
\end{equation}
where the coefficients $\alpha_n$ are complex numbers. If the basis vectors
$\ket{e_n}$ obey $\braket{e_m}{e_n} = \delta_{mn}$ the basis is called
orthonormal.

An operator $\hat{O}$ acting on the Hilbert space $\mathcal{H}$ maps vectors
onto other vectors in the Hilbert space,
\begin{equation}
  \ket{u}\in\mathcal{H} \;\longrightarrow\; \hat{O}\ket{u}\in\mathcal{H}.
\end{equation}
The corresponding adjoint operator $\hat{O}^\dagger$ is defined by
\begin{equation}
  \bra{u}\hat{O}\ket{v} = \bra{v}\hat{O}^\dagger\ket{u}^*.
\end{equation}
If the operator obeys $\hat{O}^\dagger = \hat{O}$ for all $\mathcal{H}$ it is
called self-adjoint or hermitian. An eigenvector or eigenstate $\ket{u}$ of an
operator $\hat{O}$ is a vector obeying the equation
\begin{equation}
  \hat{O}\ket{u} = \lambda\ket{u},
\end{equation}
for some complex number $\lambda$, the so-called eigenvalue. The eigenvalues of
self-adjoint operators are real and eigenvectors with different eigenvalues are
orthogonal to each other. Eigenvectors with equal eigenvalues can be made
orthogonal and after normalization the eigenvectors of a self-adjoint operator
form an orthonormal set. For a large class of operators the corresponding set
of eigenvectors provides an orthonormal basis for the Hilbert space.\footnote{We
remark that our short summary of Hilbert spaces focuses on the algebraic
properties we actually need in subsequent calculations. For more subtle aspects
such as the completeness of an eigensystem or the discussion of eigenvectors of
the continuous spectrum we refer the reader to the more mathematical literature
[5, 6].}

The unit operator $\mathbb{1}$ can be written in terms of the vectors of a
complete orthonormal basis as
\begin{equation}
  \mathbb{1} = \sum_n \ket{e_n}\bra{e_n}.
\end{equation}
This is also called the completeness relation. Finally, we define the trace
$\tr[\hat{O}]$ of an operator as
\begin{equation}
  \tr[\hat{O}] = \sum_n \bra{e_n}\hat{O}\ket{e_n},
\end{equation}
where the sum runs over the vectors of an orthonormal basis. Operators in which
the trace exists are called trace class. We stress that the trace is invariant
under a change of the basis used for its evaluation since different orthonormal
bases are related by unitary transformations.

In some of the calculations in this book we will use non-countable complete
sets of states, such as momentum eigenstates. For these sets the sums in (1.6),
(1.10), and (1.11) are replaced by integrals.

\subsection{Remarks on Hilbert spaces in particle physics}

Let us add a few remarks about the Hilbert spaces relevant to particle physics.
We do not attempt a systematic construction of these Hilbert spaces (see, e.g.,
[4]) but concentrate on discussing some properties needed here.

An important feature of Hilbert spaces in particle physics is the fact that
their vectors can be multiparticle states. The underlying physical reason is
that, when relativistic field theories are used to describe a physical system,
particle creation and annihilation processes take place. In order to
incorporate this feature, the Hilbert space is a direct sum of 0-particle,
1-particle, 2-particle, etc., Hilbert spaces. Such a Hilbert space is known as
a Fock space.

A typical example of a vector in a Fock space is a state with a pion at some
position $x$ and a proton at some other position $y$ (a 2-particle state). An
example of a 1-particle state is a state with a single $\Delta$-baryon at the
origin. States of particles may be described by their quantum numbers and their
position, but also the particle momentum can be used to characterize the state.
An example of such a case would be a state consisting of two pions with momenta
$p$ and $-p$.

The 0-particle state is referred to as the vacuum and is denoted by $\ket{0}$.
We remark that in some situations, such as spontaneous breaking of symmetries,
the vacuum state is not necessarily unique.

An important class of operators are field operators creating and annihilating
particles. In many cases the adjoint of a field operator is the field operator
for the corresponding antiparticle. An example which illustrates this feature
is the pion system: Let $O_{\pi^+}^\dagger$ be the operator that creates a
state with the quantum numbers of the $\pi^+$ meson (actually this operator is
constructed as a product of quark and antiquark field operators). The adjoint
$O_{\pi^+}$ of this operator either annihilates that state or creates a state
with the quantum numbers of its antiparticle, the $\pi^-$ meson. We will
discuss the construction of hadron operators and their properties in more
detail in Chap. 6.

We remark that all the states and operators we have discussed so far are time
independent. Their evolution in Euclidean time will be discussed now.

\subsection{Euclidean correlators}

Having prepared the ground we can now start to discuss Euclidean correlators.
Let us begin with defining the Euclidean correlator
$\ev{O_2(t)\,O_1(0)}_T$ by
\begin{equation}
  \ev{O_2(t)\,O_1(0)}_T
  = \frac{1}{Z_T}\tr\!\big[e^{-(T-t)H}O_2\,e^{-tH}O_1\big],
\end{equation}
where the normalization factor $Z_T$ is given by
\begin{equation}
  Z_T = \tr\big[e^{-TH}\big].
\end{equation}
Later we will show that in the Euclidean path integral formalism the
normalization factor $Z_T$ is a partition function of a system of statistical
mechanics. Here $O_1$ and $O_2$ can be operators that create or annihilate
states, or operators that measure observables or combinations of all of these.
The self-adjoint operator $\hat{H}$ is the Hamiltonian operator of the system
which measures the energy in the system and also governs the time evolution.
The parameters $T$ and $t$ are real, non-negative numbers denoting Euclidean
time distances of propagation. Whereas $t$ is the actual distance of interest
to us, $T$ is a formal maximal distance, which will eventually be taken to
infinity. The exponentials of the Hamiltonian operator can in most cases
(bounded operators) be defined through their power series expansion
\begin{equation}
  e^{-tH} = \sum_{j=0}^{\infty}\frac{(-t)^j}{j!}\,H^j.
\end{equation}

Let us evaluate the partition function $Z_T$. According to the definition (1.11)
of the trace we need to sandwich the operator between vectors of an orthonormal
basis and then sum over all basis vectors. The result is independent of which
particular orthonormal basis one uses. For the trace in (1.13) a natural choice
is the basis of the eigenstates of $\hat{H}$. Denoting the eigenstates of
$\hat{H}$ as $\ket{n}$, the corresponding eigenvalue equation reads
\begin{equation}
  \hat{H}\ket{n} = E_n\ket{n}.
\end{equation}
The energy eigenvalues $E_n$ are real numbers and, assuming discrete spectrum,
we order the states such that
\begin{equation}
  E_0 \leq E_1 \leq E_2 \leq E_3 \cdots
\end{equation}
The index $n$ labels all the combinations of quantum numbers describing the
states such that their energies are ordered according to (1.16).

When using the basis $\ket{n}$ for evaluating the trace in (1.13) one finds
\begin{equation}
  Z_T = \sum_n \bra{n}e^{-TH}\ket{n} = \sum_n e^{-TE_n},
\end{equation}
where in the second step we used (1.15) and the normalization of the
eigenstates, $\braket{n}{n} = 1$. The partition function $Z_T$ is simply a sum
over the exponentials of all energy eigenvalues $E_n$.

The Euclidean correlator $\ev{O_2(t)\,O_1(0)}_T$ can be evaluated in a similar
way. When computing the trace as before and inserting the unit operator in the
form (1.10) to the right of $O_2$ we obtain
\begin{align}
  \ev{O_2(t)\,O_1(0)}_T
  &= \frac{1}{Z_T}\sum_{m,n}\bra{m}e^{-(T-t)H}O_2\ket{n}\bra{n}e^{-tH}O_1\ket{m} \notag\\
  &= \frac{1}{Z_T}\sum_{m,n} e^{-(T-t)E_m}\,\bbraket{m}{O_2}{n}\,e^{-tE_n}\,\bbraket{n}{O_1}{m}.
\end{align}
The next step is to insert the expression (1.17) for $Z_T$ into the last
equation and to pull out a factor of $e^{-TE_0}$ both in the numerator and in
the denominator. This leads to
\begin{equation}
  \ev{O_2(t)\,O_1(0)}_T
  = \frac{\displaystyle\sum_{m,n}\bbraket{m}{O_2}{n}\,\bbraket{n}{O_1}{m}\,
       e^{-t\Delta E_n}\,e^{-(T-t)\Delta E_m}}
         {1 + e^{-T\Delta E_1} + e^{-T\Delta E_2} + \cdots},
\end{equation}
where we defined
\begin{equation}
  \Delta E_n = E_n - E_0.
\end{equation}
Thus, the Euclidean correlator $\ev{O_1(0)\,O_2(t)}_T$ depends only on the
energies normalized relative to the energy $E_0$ of the vacuum. It is exactly
these energy differences that can be measured in an experiment. For convenience
from now on we use $E_n$ to denote the energy differences relative to the
vacuum, i.e., we use $E_n$ instead of $\Delta E_n$. This implies that the
energy $E_0$ of the vacuum $\ket{0}$ is normalized to zero.

We now analyze (1.19) in the limit $T \to \infty$. In this limit (assuming
that the vacuum is unique and $E_1 > 0$) the denominator is equal to 1.
Similarly, in the numerator only those terms where $E_m = 0$ survive, i.e.,
only the contributions where $\ket{m} = \ket{0}$. We thus obtain
\begin{equation}
  \lim_{T\to\infty}\ev{O_2(t)\,O_1(0)}_T
  = \sum_n \bbraket{0}{O_2}{n}\,\bbraket{n}{O_1}{0}\,e^{-tE_n}.
\end{equation}
The expression which we have derived here will be central for the
interpretation of lattice field theory. It is a sum of exponentials and each
exponent corresponds to an energy level. These energies can be calculated as
follows: Let us assume that one wants to compute the energy of a proton $p$.
For $O_1$ one chooses an operator $O_p^\dagger$ which creates from the vacuum a
state with the quantum numbers of the proton and for $O_2$ the adjoint of that
operator, the operator $O_p$ that annihilates the proton. Then the matrix
element $\bra{n}O_p^\dagger\ket{0}$ vanishes for all states $\bra{n}$ that do
not have the quantum numbers of the proton. The first state $\bra{n}$ where we
find a contribution is the state describing a proton, i.e., $\bra{n} = \bra{p}$.
Further up in the ladder of states one finds excited proton states $\bra{p'}$,
$\bra{p''}$, \dots, which also have nonvanishing overlap with
$O_p^\dagger\ket{0}$. We thus obtain
\begin{equation}
  \lim_{T\to\infty}\ev{O_p(t)\,O_p(0)^\dagger}_T
  = \big|\bra{p}O_p^\dagger\ket{0}\big|^2 e^{-tE_p}
    + \big|\bra{p'}O_p^\dagger\ket{0}\big|^2 e^{-tE_{p'}} + \cdots,
\end{equation}
where we made use of (1.8) to simplify the coefficients in front of the
exponentials. For sufficiently large $t$ the sub-leading terms are strongly
suppressed since $E_{p'} > E_p$. Therefore, at sufficiently large $t$ we obtain
$E_p$ from the exponential decay of the Euclidean correlator. When using other
operators $O_1$, $O_2$ that create other states from the vacuum, we can extract
the ground state energies for all particles in the spectrum of the theory in a
way similar to what we have just demonstrated for the proton. Furthermore, one
can also study the matrix elements $\bra{n}\hat{O}\ket{0}$. Since $\hat{O}$ can
be a product of several operators, a wide range of matrix elements is
accessible. Such techniques will be discussed in great detail in Chaps. 6
and 11.

This completes our discussion of the properties of Euclidean correlators. In
the next section we introduce the path integral as a technique for calculating
the correlators (1.12).

Let us conclude this section with a short remark on the relation of our
Euclidean correlators to real time, Minkowski quantum mechanics. In the
Heisenberg picture, where the operators are time dependent, an operator at time
$\tau$ is given by (note that we use $\hbar \equiv 1$)
\begin{equation}
  \hat{O}(\tau) = e^{i\tau H}\,\hat{O}\,e^{-i\tau H}.
\end{equation}
Thus, the Euclidean correlator (1.12) can be seen as the correlator of operator
$O_1$ at time 0 with operator $O_2$ at imaginary time $t = i\tau$ times an
extra time transporter $\exp(-TH)$ which projects to the vacuum for
$T \to \infty$. Four-vectors with real-time $\tau$ correspond to Minkowski
metric. When one switches to imaginary time $t = i\tau$, the relative minus
sign between time- and space-components vanishes and the metric becomes
Euclidean; thus, the name Euclidean correlators for the objects defined in
(1.12).

In the next section, where we discuss the path integral representation of the
Euclidean correlator, we will encounter a similar transformation of the metric.
The change from real to imaginary time is often referred to as Wick rotation.
We stress, however, that it is not essential to interpret $t$ as time at all.
The expression (1.12) can be simply viewed as a convenient mathematical tool
which allows one to extract energy levels and matrix elements of operators.

\section{The path integral for a quantum mechanical system}

The Euclidean correlators (1.12) which we have introduced and analyzed in
Hilbert space in the last section are now rewritten as path integrals -- a
concept introduced by Feynman. In Sect. 1.3 we will derive the second key
equation of lattice field theory, (1.2). However, as a preparatory step, in
this section we first derive the path integral for a simple quantum mechanical
system. This allows us to focus on the essential steps of the construction
without too much technicalities. In Sect. 1.3 we will repeat the derivation of
the path integral for a scalar field theory.

The quantum mechanical system we use in our initial presentation of the path
integral is a single particle in a potential $U$. To simplify things further we
restrict the motion of the particle to the $x$-axis only. The system is
described by the Hamiltonian
\begin{equation}
  \hat{H} = \frac{\hat{p}^2}{2m} + U.
\end{equation}
The canonical quantization condition is given by the commutator
\begin{equation}
  [\hat{x}, \hat{p}] = i,
\end{equation}
(we use $\hbar \equiv 1$) and implies that the momentum operator is given by
\begin{equation}
  \hat{p} = -i\frac{d}{dx}.
\end{equation}

For this system we want to compute the partition function according to (1.13)
\begin{equation}
  Z_T = \tr\big[e^{-TH}\big] = \int dx\,\bra{x}e^{-TH}\ket{x}.
\end{equation}
The second step in this equation already expresses our intent to use position
states $\ket{x}$ to compute the trace.

Let us begin with the evaluation of the necessary matrix elements for the free
case:
\begin{equation}
  \bra{x}e^{-t\hat{p}^2/2m}\ket{y}
  = \int dp\,\bbraket{x}{e^{-t\hat{p}^2/2m}}{p}\braket{p}{y}
  = \int dp\,\braket{x}{p}\,\braket{p}{y}\,e^{-tp^2/2m}.
\end{equation}
In the first step we have inserted the unit operator according to (1.10) using
the momentum eigenstates $\ket{p}$. In the second step we have used
$\hat{p}\ket{p} = p\ket{p}$. Note that in the exponent on the right-hand side
of (1.28), $p$ is no longer an operator but a number -- the momentum.

Using the real space representation of the momentum eigenstates, i.e., the
plane waves
\begin{equation}
  \braket{x}{p} = \frac{1}{\sqrt{2\pi}}\,e^{ipx},
\end{equation}
we find for our matrix element
\begin{equation}
  \bra{x}e^{-t\hat{p}^2/2m}\ket{y}
  = \frac{1}{2\pi}\int dp\,e^{ip(x-y)}e^{-tp^2/2m}
  = \sqrt{\frac{m}{2\pi t}}\,e^{-(x-y)^2m/(2t)}.
\end{equation}
In the last step we used the well-known Gaussian integral ($c$, $b$ real,
$c > 0$)
\begin{equation}
  \int_{-\infty}^{\infty}dx\,e^{-cx^2}\,e^{\pm ibx}
  = \sqrt{\frac{\pi}{c}}\,e^{-b^2/4c}.
\end{equation}

Let us now look at the general case with nonvanishing potential $U$. The
problem here is that the two terms in the Hamiltonian (1.24) do not commute
with each other, making the evaluation of the exponential of $\hat{H}$
somewhat involved. We split the potential $U$ symmetrically, and for an
infinitesimal Euclidean timestep $\varepsilon$ we write
\begin{equation}
  e^{-\varepsilon H} = e^{-\varepsilon U/2}\,e^{-\varepsilon\hat{p}^2/2m}\,
  e^{-\varepsilon U/2}\,\big(1 + \order{\varepsilon}\big).
\end{equation}
This formula can be obtained from expanding both sides in $\varepsilon$. We
abbreviate the leading term as
\begin{equation}
  W_\varepsilon = e^{-\varepsilon U/2}\,e^{-\varepsilon\hat{p}^2/2m}\,
  e^{-\varepsilon U/2},
\end{equation}
and find
\begin{equation}
  \begin{split}
    \bra{x}W_\varepsilon\ket{y}
    &= e^{-\varepsilon U(x)/2}\,\bra{x}e^{-\varepsilon\hat{p}^2/2m}\ket{y}\,
       e^{-\varepsilon U(y)/2}\\
    &= \sqrt{\frac{m}{2\pi\varepsilon}}\,
       e^{-\varepsilon U(x)/2}\,e^{-\varepsilon U(y)/2}\,e^{-(x-y)^2m/(2\varepsilon)}.
  \end{split}
\end{equation}
In the first step we have used $\hat{U}\ket{x} = U(x)\ket{x}$, where $U(x)$ is
a real number -- the value of the potential at $x$. In the second step we
inserted (1.30).

We can build up a finite Euclidean timestep $T$ from infinitesimal steps
$\varepsilon$ using the Trotter formula (see, e.g., [7] for a proof),
\begin{equation}
  e^{-TH} = \lim_{N_T\to\infty} W_\varepsilon^{N_T}
  \qquad \text{with} \quad \varepsilon = \frac{T}{N_T}.
\end{equation}

We obtain
\begin{align}
  Z_T &= \int dx_0\,\bra{x_0}e^{-TH}\ket{x_0}
        = \lim_{N_T\to\infty}\int dx_0\,\bra{x_0}W_\varepsilon^{N_T}\ket{x_0} \notag\\
      &= \lim_{N_T\to\infty}\int dx_0\cdots dx_{N_T-1}\,
         \bra{x_0}W_\varepsilon\ket{x_1}\cdots \notag\\
      &\qquad \bra{x_{N_T-2}}W_\varepsilon\ket{x_{N_T-1}}
         \bra{x_{N_T-1}}W_\varepsilon\ket{x_0} \notag\\
      &= \lim_{N_T\to\infty} C^{N_T}\int dx_0\cdots dx_{N_T-1}\, \notag\\
      &\quad \exp\!\Big(-\varepsilon\sum_{j=0}^{N_T-1}
           \Big[\frac{m}{2}\Big(\frac{x_j - x_{j+1}}{\varepsilon}\Big)^2 + U(x_j)\Big]\Big).
\end{align}
In this derivation we have inserted unit operators in the form
$\int dx_i \ket{x_i}\bra{x_i}$ and then used the result (1.34) with the
abbreviation $C = \sqrt{m/(2\pi\varepsilon)}$.

Our expression (1.36) contains the limit $N_T \to \infty$. Since we are aiming
at a numerical evaluation of our path integral with a computer, the number of
steps $N_T$ we can use has to be finite. Hence we define an approximation of
the partition function (we periodically identify $x_{N_T} \equiv x_0$)
\begin{equation}
  Z_T^\varepsilon = C^{N_T}\int dx_0\cdots dx_{N_T-1}\,
  \exp\!\Big(-\varepsilon\sum_{j=0}^{N_T-1}
    \Big[\frac{m}{2}\Big(\frac{x_{j+1} - x_j}{\varepsilon}\Big)^2 + U(x_j)\Big]\Big),
\end{equation}
where the stepsize $\varepsilon$ and the Euclidean time are related by
$T = \varepsilon N_T$.

The sum in the exponent has an interesting interpretation. Obviously, for
smooth paths
\begin{align}
  \frac{x_{j+1} - x_j}{\varepsilon} &= \dot{x}(t) + \order{\varepsilon}
    \qquad \text{with} \quad t = j\varepsilon, \notag\\
  \sum_{j=0}^{N_T-1}\varepsilon\;\cdots\; &= \int_0^T dt\;\cdots\;
    + \order{\varepsilon} \qquad \text{with} \quad T = N_T\varepsilon.
\end{align}
Thus we can identify
\begin{equation}
  \sum_{j=0}^{N_T-1}\varepsilon
  \Big[\frac{m}{2}\Big(\frac{x_{j+1} - x_j}{\varepsilon}\Big)^2 + U(x_j)\Big]
  = \int_0^T dt\;\Big[\frac{m}{2}\dot{x}(t)^2 + U(x(t))\Big] + \order{\varepsilon}.
\end{equation}
The expression on the right-hand side is called the Euclidean action $S_E$. It
is obtained from the usual action $S$ when switching from real-time $\tau$ to
imaginary-time $t = i\tau$ and rotating the contour of integration:
\begin{equation}
  \begin{split}
    S[x,\dot{x}] &= \int_0^T d\tau\,
      \Big[\frac{m}{2}\dot{x}(\tau)^2 - U(x(\tau))\Big]\\
    &\;\to\; i\int_0^T dt\,
      \Big[\frac{m}{2}\dot{x}(t)^2 + U(x(t))\Big] = i\,S_E[x,\dot{x}].
  \end{split}
\end{equation}

\begin{figure}[htbp]
  \centering
  \includegraphics[width=0.8\textwidth]{fig11.png}
  \caption*{Fig. 1.1. \emph{A discretized path contributing in (1.37)}.}
  \label{fig:11}
\end{figure}

Let us summarize: We have found that we can construct an approximation
$Z_T^\varepsilon$ of the partition function in the form of (1.37). In this
approximation we divide the Euclidean time interval $T$ into $N_T$ steps of
size $\varepsilon$ (using periodic boundary conditions). At each step we insert
a variable $x_j$ which we integrate from $-\infty$ to $+\infty$. The collection
of values $x_j$ can be interpreted as a path (compare Fig. 1.1) and the
integral is over all possible paths. The integrand is the exponential (1.37)
of the Euclidean action for the discretized path.

In the next section we will develop a path integral representation for a system
much closer to our target theory QCD. We will consider a scalar field theory.
The technical steps of the construction -- free case, infinitesimal timesteps,
Trotter formula -- will be the same. However, since we are dealing with a
quantum field theory with many degrees of freedom, the calculation, in
particular the corresponding notation, will be a little bit more involved.

\section{The path integral for a scalar field theory}

We now derive the key equation (1.2) for a scalar field theory. Although this
is already a field theory, it has a simpler structure than QCD. This allows us
to focus on the central points of the derivation of the path integral for a
field theory without having to worry too much about technicalities. Although
our presentation is self contained we recommend [8, 9] for an introductory
reading about the canonical quantization of the scalar field, which is our
starting point.

\subsection{The Klein--Gordon field}

Let us begin by recalling the classical action and equations of motion of the
scalar field in Minkowski space. We consider a real scalar field
$\Phi(t,x)$. The corresponding action $S$ is an integral over space-time
\begin{equation}
  S = \int dt\,d^3x\;\mathcal{L}\big(\Phi(t,x), \partial_\mu\Phi(t,x)\big),
\end{equation}
with the Lagrangian density $\mathcal{L}$ given by
\begin{align}
  \mathcal{L}(\Phi, \partial_\mu\Phi)
  &= \frac{1}{2}(\partial_\mu\Phi)(\partial^\mu\Phi) - \frac{m^2}{2}\Phi^2 - V(\Phi) \notag\\
  &= \frac{1}{2}\dot{\Phi}^2 - \frac{1}{2}(\nabla\Phi)^2 - \frac{m^2}{2}\Phi^2 - V(\Phi).
\end{align}
This is a system of coupled oscillators. We stress again that this expression
is in Minkowski space, i.e., in the second step we used the metric
$g_{\mu\nu} = \mathrm{diag}(1,-1,-1,-1)$ to rewrite the kinetic term. In the action we
also allow for a potential term $V(\Phi)$. A standard example is
$V(\Phi(t,x)) = \lambda\,\Phi(t,x)^4$.

Using the Euler--Lagrange equations
\begin{equation}
  \partial_\mu\frac{\partial\mathcal{L}}{\partial(\partial_\mu\Phi)}
  - \frac{\partial\mathcal{L}}{\partial\Phi} = 0,
\end{equation}
we can derive the classical equation of motion
\begin{equation}
  \partial_\mu(\partial^\mu\Phi) + m^2\Phi + V'(\Phi) = 0.
\end{equation}
In the absence of a potential this is the Klein--Gordon equation.

As a next step we quantize the system using the canonical formalism. For this
purpose we need the Hamiltonian which we obtain from the Lagrangian density.
The canonical momentum $\Pi(t,x)$ is defined by
\begin{equation}
  \Pi(t,x) = \frac{\partial}{\partial\dot{\Phi}(t,x)}
    \mathcal{L}\big(\Phi(t,x), \partial_\mu\Phi(t,x)\big) = \dot{\Phi}(t,x),
\end{equation}
and the Hamiltonian function $H$ is obtained as the Legendre transform
\begin{align}
  H &= \int d^3x\,\Pi(t,x)\,\dot{\Phi}(t,x)
      - \int d^3x\;\mathcal{L}\big(\Phi(t,x), \partial_\mu\Phi(t,x)\big) \notag\\
    &= \int d^3x\;\Big[\frac{1}{2}\Pi(t,x)^2
       + \frac{1}{2}\big(\nabla\Phi(t,x)\big)^2
       + \frac{m^2}{2}\Phi(t,x)^2 + V\big(\Phi(t,x)\big)\Big].
\end{align}
The Hamiltonian function is the starting point for the quantization of the
system in the continuum. Upon quantization, the Hamiltonian function $H$ and
the fields $\Phi(t,x)$ and $\Pi(t,x)$ turn into Schr\"odinger operators
$\hat{H}$, $\hat{\Phi}(x)$, $\hat{\Pi}(x)$. The Hamiltonian operator reads (up
to normal ordering)
\begin{equation}
  \hat{H} = \int d^3x\;\Big[\frac{1}{2}\hat{\Pi}(x)^2
    + \frac{1}{2}\big(\nabla\hat{\Phi}(x)\big)^2
    + \frac{m^2}{2}\hat{\Phi}(x)^2 + V\big(\hat{\Phi}(x)\big)\Big].
\end{equation}
All time arguments are gone, and the time evolution of operators is given by
(1.23). The operators $\hat{\Phi}$, $\hat{\Pi}$ obey the canonical equal time
commutation relations ($\hbar \equiv 1$)
\begin{equation}
  \big[\hat{\Phi}(x), \hat{\Pi}(y)\big] = i\,\delta(x-y), \qquad
  \big[\hat{\Phi}(x), \hat{\Phi}(y)\big] = 0, \qquad
  \big[\hat{\Pi}(x), \hat{\Pi}(y)\big] = 0.
\end{equation}
We now leave the canonical path and follow a different strategy by introducing
the lattice as a cutoff before we quantize the system.

\subsection{Lattice regularization of the Klein--Gordon Hamiltonian}

In order to calculate something from a quantum field theory in a mathematically
well-defined way, an ultraviolet regulator must be introduced. Such a regulator
is necessary to make the expressions finite. In continuum perturbation theory
the regulator can, e.g., be introduced by using dimensional regularization,
Pauli--Villars regularization or a momentum cutoff. A different approach is
chosen in lattice field theory.

The central idea is to replace the continuous space by a 3D finite lattice
$\Lambda^3$:
\begin{equation}
  x \;\Rightarrow\; an, \qquad n_i = 0, 1, \ldots, N-1
  \quad \text{for} \quad i = 1, 2, 3.
\end{equation}
The lattice constant $a$ has the physical dimension of length. The vector $n$
with integer-valued components, $n_1$, $n_2$, $n_3$, labels the lattice sites.
The operators $\hat{\Phi}(n)$ and $\hat{\Pi}(n)$ now live only on the lattice
sites and we stress this fact by replacing the space argument $x$ by the label
of the lattice site $n$. Since each component $n_i$ runs from 0 to $N-1$, we
have a total of $N^3$ lattice sites and the lattice system thus has $2N^3$
degrees of freedom ($\hat{\Phi}$ and $\hat{\Pi}$ on each lattice site). For any
finite system one has to decide what to do at the boundary. We choose periodic
boundary conditions, i.e., we identify $n_j = N$ with $n_j = 0$.

For a lattice regularization of the Hamiltonian operator (1.47) we also need
to discretize the derivatives $\nabla\hat{\Phi}(x)$. Using the Taylor series
expansion of $\hat{\Phi}(x)$, it is easy to see that for small lattice constant
$a$ the derivative may be approximated as
\begin{equation}
  \partial_j\hat{\Phi}(x)
  = \frac{\hat{\Phi}(n+\hat{\jmath}) - \hat{\Phi}(n-\hat{\jmath})}{2a}
    + \order{a^2},
\end{equation}
where $\hat{\jmath}$ denotes the unit vector in the $j$-direction. Using this
definition for the derivative, we obtain a lattice version of the Hamiltonian
operator (1.47):
\begin{equation}
  \hat{H} = a^3\sum_{n\in\Lambda^3}\Big[
    \frac{1}{2}\hat{\Pi}(n)^2
    + \frac{1}{2}\sum_{j=1}^3
      \Big(\frac{\hat{\Phi}(n+\hat{\jmath}) - \hat{\Phi}(n-\hat{\jmath})}{2a}\Big)^2
    + \frac{m^2}{2}\hat{\Phi}(n)^2 + V\big(\hat{\Phi}(n)\big)\Big].
\end{equation}
The integral $\int d^3x$ has been replaced by the sum
$a^3\sum_{n\in\Lambda^3}$. The Hamiltonian describes a system with a finite
number of degrees of freedom $\hat{\Phi}(n)$ and their corresponding conjugate
momenta $\hat{\Pi}(n)$. The canonical quantization condition is given by the
commutators
\begin{equation}
  \big[\hat{\Phi}(n), \hat{\Pi}(m)\big] = i\,a^{-3}\delta_{nm},
\end{equation}
\begin{equation}
  \big[\hat{\Phi}(n), \hat{\Phi}(m)\big] = 0, \qquad
  \big[\hat{\Pi}(n), \hat{\Pi}(m)\big] = 0.
\end{equation}
When comparing the new quantization conditions (1.52) and (1.53) to their
continuum counterparts one sees that the Dirac-delta in (1.48) is replaced by
Kronecker deltas $\delta_{nm} = \delta_{n_1 m_1}\delta_{n_2 m_2}\delta_{n_3 m_3}$,
as is usual for discrete degrees of freedom. The factor $a^{-3}$ on the
right-hand side of (1.52) is needed for getting the dimensions right. When
denoting the dimension of length by $\ell$ we find for the dimensions
$[\,\cdot\,]$ of the quantities involved (using $\hbar \equiv 1$):
\begin{equation}
  [S] = 1, \qquad [H] = \frac{1}{\ell}, \qquad [a] = \ell, \qquad
  [\hat{\Phi}] = \frac{1}{\ell}, \qquad [\hat{\Pi}] = \frac{1}{\ell^2}, \qquad
  [m] = \frac{1}{\ell}.
\end{equation}
Thus the factor $a^{-3}$ in (1.52) is necessary in order to have objects of
equal dimensions on both sides. This factor also matches the dimensions in the
continuum formula (1.48) since the Dirac-delta there also has dimension
$\ell^{-3}$.

The commutators (1.52) and (1.53) are obeyed when the momentum operators are
represented as derivatives
\begin{equation}
  \hat{\Pi}(n) = -\frac{i}{a^3}\frac{\partial}{\partial\Phi(n)}.
\end{equation}

Using this representation we can write the Hamiltonian operator for the
Klein--Gordon field on the lattice in its final form as
\begin{align}
  \hat{H} &= \hat{H}_0 + \hat{U}, \notag\\
  \hat{H}_0 &= a^3\sum_{n\in\Lambda^3}\frac{1}{2}
    \Big(-\frac{i}{a^3}\frac{\partial}{\partial\Phi(n)}\Big)^2
    = -\sum_{n\in\Lambda^3}\frac{1}{2a^3}\frac{\partial^2}{\partial\Phi(n)^2}, \notag\\
  \hat{U} &= a^3\sum_{n\in\Lambda^3}\Big[
    \frac{1}{2}\sum_{j=1}^3
      \Big(\frac{\hat{\Phi}(n+\hat{\jmath}) - \hat{\Phi}(n-\hat{\jmath})}{2a}\Big)^2
    + \frac{m^2}{2}\hat{\Phi}(n)^2 + V\big(\hat{\Phi}(n)\big)\Big].
\end{align}
We have split the Hamilton operator into two parts: the free part $\hat{H}_0$,
which contains the derivatives, and the interaction part $\hat{U}$. This
splitting is useful for the calculation of the Euclidean correlators.

It is convenient to introduce a set of eigenstates of the field operators
$\hat{\Phi}(n)$, equivalent to the position states $\ket{x}$ of quantum
mechanics. The condition (1.52) is a generalization of the condition
$[\hat{x}, \hat{p}] = i\hbar$ in ordinary quantum mechanics. In a similar way
the states $\ket{\Phi}$ are a generalization of the position states $\ket{x}$.
The action of the operators $\hat{\Phi}(n)$ on $\ket{\Phi}$ is defined as
\begin{equation}
  \hat{\Phi}(n)\ket{\Phi} = \Phi(n)\ket{\Phi}.
\end{equation}
The application of the operator $\hat{\Phi}(n)$ on an eigenstate $\ket{\Phi}$
thus gives as the eigenvalue the value $\Phi(n)$ of the field at the lattice
point $n$. Consequently, the state $\ket{\Phi}$ is labeled by the values of the
field at all lattice points, i.e., by the set $\{\Phi(n), n \in \Lambda^3\}$.
The states $\ket{\Phi}$ are orthogonal and complete:
\begin{equation}
  \braket{\Phi'}{\Phi} = \delta(\Phi' - \Phi)
  \equiv \prod_{n\in\Lambda^3}\delta\big(\Phi'(n) - \Phi(n)\big),
\end{equation}
\begin{equation}
  \mathbb{1} = \int_{-\infty}^{\infty}D\Phi\;\ket{\Phi}\bra{\Phi}
  \qquad \text{with} \quad D\Phi = \prod_{n\in\Lambda^3}d\Phi(n).
\end{equation}

\subsection{The Euclidean time transporter for the free case}

Having discussed the Hamiltonian of the scalar field and its lattice
discretization, we can now return to our original problem, the calculation of
the Euclidean correlators (1.12) and the partition function (1.13). In both
these expressions we encounter Euclidean time transporters $\exp(-tH)$. For the
full Hamiltonian operator $\hat{H}$ the calculation of matrix elements of the
time transporter is not possible in closed form. However, for the free part
$\hat{H}_0$ this is an easy exercise which we perform using the eigenstates of
$\hat{H}_0$.

The free Hamiltonian operator $\hat{H}_0$, as represented in (1.56), is a sum
over second derivatives, i.e., a sum over squares of momentum operators. It is
well-known from quantum mechanics that the eigenstates of a derivative operator
are the momentum eigenstates $\ket{\Pi}$: plane waves.\footnote{These are not
plane waves in real space, but plane waves of the fields $\Phi(n)$.} Since
$\hat{H}_0$ is a sum over $N^3$ derivatives, the eigenstate is a product over
$N^3$ plane waves, one for each field component $\Phi(n)$. Each plane wave is
characterized by a wave number $\Pi(n)$. The whole eigenstate is then labeled
by the set of all these wave numbers and we denote the eigenstate as
$\ket{\Pi}$, which in field representation is given by
\begin{equation}
  \braket{\Phi}{\Pi}
  = \prod_{n\in\Lambda^3}\sqrt{\frac{a^3}{2\pi}}\;e^{i a^3\Pi(n)\Phi(n)}.
\end{equation}
The eigenstates in (1.60) are normalized such that the unit operator reads
\begin{equation}
  \mathbb{1} = \int_{-\infty}^{\infty}D\Pi\;\ket{\Pi}\bra{\Pi}
  \qquad \text{with} \quad D\Pi = \prod_{n\in\Lambda^3}d\Pi(n).
\end{equation}

The action of the free Hamiltonian $\hat{H}_0$ on $\ket{\Pi}$ gives rise to
\begin{equation}
  \bra{\Phi}\hat{H}_0\ket{\Pi}
  = -\frac{1}{2a^3}\sum_{n\in\Lambda^3}
    \frac{\partial^2}{\partial\Phi(n)^2}\braket{\Phi}{\Pi}
  = \frac{a^3}{2}\sum_{n\in\Lambda^3}\Pi(n)^2\braket{\Phi}{\Pi}.
\end{equation}
Thus the eigenvalues of $\hat{H}_0$ are given by
$(a^3/2)\sum_n\Pi(n)^2$. It is important to realize that in both Eqs. (1.60)
and (1.62) the quantities $\Phi(n)$ and $\Pi(n)$ are no longer operators, but
real numbers. They are the values of the field and the wave number at the
lattice site $n$.

With the help of the eigenstates (1.60) we can now compute matrix elements of
the Euclidean time transporter of the free theory
\begin{align}
  \bra{\Phi'}e^{-t\hat{H}_0}\ket{\Phi}
  &= \int D\Pi\;\bra{\Phi'}e^{-t\hat{H}_0}\ket{\Pi}\braket{\Pi}{\Phi} \notag\\
  &= \int D\Pi\;\braket{\Phi'}{\Pi}\,\braket{\Pi}{\Phi}\,
     e^{-t\frac{a^3}{2}\sum_n \Pi(n)^2},
\end{align}
where in the first step we have inserted the unit operator written in the form
(1.61) and in the second step the eigenvalue equation (1.62) was used.
Inserting (1.60) into the last expression, we have a product of Gaussian
integrals that can be solved:
\begin{align}
  \bra{\Phi'}e^{-t\hat{H}_0}\ket{\Phi}
  &= \prod_{n\in\Lambda^3}\frac{a^3}{2\pi}
     \int d\Pi(n)\;e^{i a^3\Pi(n)(\Phi'(n)-\Phi(n))}
     e^{-t a^3\Pi(n)^2/2} \notag\\
  &= \prod_{n\in\Lambda^3}\sqrt{\frac{a^3}{2\pi t}}\;
     e^{-a^3/(2t)(\Phi'(n)-\Phi(n))^2} \notag\\
  &= \Big(\frac{a^3}{2\pi t}\Big)^{N^3/2}
     e^{-a^3/(2t)\sum_n(\Phi'(n)-\Phi(n))^2},
\end{align}
where in the last step we used the Gaussian integral (1.31).

\subsection{Treating the interaction term with the Trotter formula}

When we switch from the free case to including the potential $\hat{U}$ we again
have to build up a finite step $t$ in Euclidean time from infinitesimal steps
using the Trotter formula (1.35). Like for the quantum mechanical system of the
last section we define
\begin{equation}
  W_\varepsilon = e^{-\varepsilon\hat{U}/2}\,e^{-\varepsilon\hat{H}_0}\,
  e^{-\varepsilon\hat{U}/2}.
\end{equation}
Using the Trotter formula (1.35) we find
\begin{align}
  \bra{\Phi'}e^{-t\hat{H}}\ket{\Phi}
  &= \lim_{n_t\to\infty}\bra{\Phi'}W_\varepsilon^{n_t}\ket{\Phi} \notag\\
  &= \lim_{n_t\to\infty}\int D\Phi_1 D\Phi_2\cdots D\Phi_{n_t-1}\,
     \bra{\Phi'}W_\varepsilon\ket{\Phi_{n_t-1}} \notag\\
  &\qquad \bra{\Phi_{n_t-1}}W_\varepsilon\ket{\Phi_{n_t-2}}\cdots
     \bra{\Phi_1}W_\varepsilon\ket{\Phi}.
\end{align}
In last step we have inserted the unit operator $n_t - 1$ times in the form
(1.59).

What remains to be done is the calculation of the matrix elements appearing
in (1.66), i.e., $\bra{\Phi_{i+1}}W_\varepsilon\ket{\Phi_i}$. We first note
that, due to (1.56) and (1.57),
\begin{equation}
  \hat{U}\ket{\Phi} = U[\Phi]\ket{\Phi},
\end{equation}
with
\begin{equation}
  U[\Phi] = a^3\sum_{n\in\Lambda^3}\Big[
    \frac{1}{2}\sum_{j=1}^3
      \Big(\frac{\Phi(n+\hat{\jmath}) - \Phi(n-\hat{\jmath})}{2a}\Big)^2
    + \frac{m^2}{2}\Phi(n)^2 + V(\Phi(n))\Big].
\end{equation}
We stress that $U[\Phi]$ is a number, i.e., the classical interaction term
evaluated for the classical fields $\Phi(n)$. Using (1.64) and (1.67) we obtain
\begin{align}
  \bra{\Phi_{i+1}}W_\varepsilon\ket{\Phi_i}
  &= \bra{\Phi_{i+1}}e^{-\varepsilon\hat{U}/2}\,e^{-\varepsilon\hat{H}_0}\,
     e^{-\varepsilon\hat{U}/2}\ket{\Phi_i} \notag\\
  &= e^{-\varepsilon U[\Phi_{i+1}]/2}\,
     \bra{\Phi_{i+1}}e^{-\varepsilon\hat{H}_0}\ket{\Phi_i}\,
     e^{-\varepsilon U[\Phi_i]/2} \notag\\
  &= C^N\,e^{-\varepsilon U[\Phi_i]/2}\,
     e^{-a^3/(2\varepsilon)\sum_n(\Phi(n)_i - \Phi(n)_{i+1})^2}\,
     e^{-\varepsilon U[\Phi_{i+1}]/2}.
\end{align}
We have introduced the abbreviation $C$ for the factor $\sqrt{a^3/(2\pi
\varepsilon)}$. Again we emphasize that the right-hand side of (1.69) is a real
number. It is the matrix element of the Euclidean time transporter for
infinitesimal time $W_\varepsilon$. These matrix elements may be combined in
(1.66) to obtain the matrix element for a transport by a finite amount of
Euclidean time.

\subsection{Path integral representation for the partition function}

Let us now use the two Eqs. (1.66) and (1.69) to compute the partition function
$Z_T$ defined in (1.13). When the trace in this equation is evaluated using
states $\ket{\Phi_0}$, we find
\begin{equation}
  Z_T = \int D\Phi_0\;\bra{\Phi_0}e^{-TH}\ket{\Phi_0}
  = \lim_{N_T\to\infty}\int D\Phi_0\;\bra{\Phi_0}W_\varepsilon^{N_T}\ket{\Phi_0},
\end{equation}
with $T = N_T\varepsilon$. We have already remarked that we are aiming at
implementing the path integral in a numerical simulation and thus we cannot
perform the limit $N_T \to \infty$, since a computer can only cope with a
finite number of degrees of freedom. Therefore, we have to work with a finite
resolution $\varepsilon$ and thus obtain only an approximation of the exact
partition function. We denote this approximation by $Z_T^\varepsilon$. Inserting
again $N_T - 1$ sets of complete states and using the expression (1.69) for the
matrix elements $\bra{\Phi_{i+1}}W_\varepsilon\ket{\Phi_i}$ we obtain
\begin{align}
  Z_T^\varepsilon &= \int D\Phi_0\;\bra{\Phi_0}W_\varepsilon^{N_T}\ket{\Phi_0} \notag\\
  &= \int D\Phi_0\cdots D\Phi_{N_T-1}\,
     \bra{\Phi_0}W_\varepsilon\ket{\Phi_{N_T-1}}
     \bra{\Phi_{N_T-1}}W_\varepsilon\ket{\Phi_{N_T-2}}\cdots
     \bra{\Phi_1}W_\varepsilon\ket{\Phi_0} \notag\\
  &= C^{N^3N_T}\int D\Phi_0\cdots D\Phi_{N_T-1}\;e^{-S_E[\Phi]},
\end{align}
with
\begin{equation}
  S_E[\Phi] = \frac{a^3}{2\varepsilon}\sum_{n\in\Lambda^3}
    \sum_{j=0}^{N_T-1}\big(\Phi(n)_{j+1} - \Phi(n)_j\big)^2
    + \varepsilon\sum_{j=0}^{N_T-1}U[\Phi_j].
\end{equation}
We remark that in the first sum we have continued the index $j$ periodically by
identifying $j = N_T$ with $j = 0$.

The expression for $S_E$ in (1.72) can be simplified by introducing a more
compact notation $\Phi(n, n_4) \equiv \Phi(n)_{n_4}$. This corresponds to
introducing a 4D lattice $\Lambda$ defined as
\begin{equation}
  \Lambda = \big\{(n, n_4)\;\big|\; n\in\Lambda^3,\;
    n_4 = 0, 1, \ldots, N_T - 1\big\},
\end{equation}
where we have also relabeled $j$ to $n_4$. Note that the lattice $\Lambda$ is
periodic in all four directions since it inherits the spatial periodicity of
the spatial components of $\Lambda^3$. With this notation $S_E$ reads
\begin{equation}
  \begin{split}
    S_E[\Phi] = \varepsilon a^3\sum_{(n,n_4)\in\Lambda}\Big[
    &\frac{1}{2}\Big(\frac{\Phi(n,n_4+1) - \Phi(n,n_4)}{\varepsilon}\Big)^2\\
    &+ \frac{1}{2}\sum_{j=1}^3
      \Big(\frac{\Phi(n+\hat{\jmath},n_4) - \Phi(n-\hat{\jmath},n_4)}{2a}\Big)^2\\
    &+ \frac{m^2}{2}\Phi(n,n_4)^2 + V\big(\Phi(n,n_4)\big)\Big].
  \end{split}
\end{equation}
As in the quantum mechanical problem of the last section, the expression (1.74)
for $S_E$ again is a discretization of the Euclidean action: The first term on
the right-hand side can be interpreted as the square of a discretized
derivative in 4-direction. The sum over $\Lambda$ together with the factor
$\varepsilon a^3$ is a discretization of an integral over space and Euclidean
time. With these two identifications the expression (1.74) is very similar to
the action of the Klein--Gordon field given in (1.41) and (1.42). The only
differences are an overall minus sign and a relative sign between the time
derivative in (1.42) and the rest of the terms. If one rotates the time in
(1.42) into the imaginary axis, i.e., one replaces $t$ by $i\tau$, then this
relative minus sign is gone. Thus $S_E[\Phi]$ is the action of the system for
Euclidean time -- the Euclidean action.

Also the measure in (1.71) can be written in simpler form as
\begin{equation}
  D[\Phi] = \prod_{(n,n_4)\in\Lambda}d\Phi(n,n_4),
\end{equation}
and the expression for the partition function turns into
\begin{equation}
  Z_T^\varepsilon = C^{N^3N_T}\int D[\Phi]\;e^{-S_E[\Phi]}.
\end{equation}

\subsection{Including operators in the path integral}

We may evaluate the right-hand side of (1.2) analogously to the method used for
the partition function $Z_T$. In particular, we again use the Trotter formula
to write the numerator of the Euclidean correlator as
\begin{equation}
  \tr\big[e^{-(T-t)H}O_2\,e^{-tH}O_1\big]
  = \lim_{N_T\to\infty}\tr\big[W_\varepsilon^{N_T-n_t}O_2\,W_\varepsilon^{n_t}O_1\big],
\end{equation}
where the number of steps $N_T$ and $n_t$ are related to $T$ and $t$ by
$T = \varepsilon N_T$ and $t = \varepsilon n_t$. Similar to the approximation
$Z_T^\varepsilon$ we also introduce a Euclidean correlator with a finite
temporal resolution $\varepsilon$. It reads
\begin{align}
  \ev{O_2(t)\,O_1(0)}_T^\varepsilon
  &= \frac{1}{Z_T^\varepsilon}\tr\big[W_\varepsilon^{N_T-n_t}O_2\,W_\varepsilon^{n_t}O_1\big] \notag\\
  &= \frac{1}{Z_T^\varepsilon}\int D\Phi_0\cdots D\Phi_{N_T-1}
     D\tilde{\Phi}_0\,D\tilde{\Phi}_{n_t}\, \notag\\
  &\qquad \bra{\Phi_0}W_\varepsilon\ket{\Phi_{N_T-1}}
     \bra{\Phi_{N_T-1}}W_\varepsilon\ket{\Phi_{N_T-2}}\cdots \notag\\
  &\qquad \times \bra{\Phi_{n_t+1}}W_\varepsilon\ket{\tilde{\Phi}_{n_t}}\,
     \bra{\tilde{\Phi}_{n_t}}O_2\ket{\Phi_{n_t}} \notag\\
  &\qquad\quad \bra{\Phi_{n_t}}W_\varepsilon\ket{\Phi_{n_t-1}}\cdots
     \bra{\Phi_1}W_\varepsilon\ket{\tilde{\Phi}_0}\,
     \bra{\tilde{\Phi}_0}O_1\ket{\Phi_0}.
\end{align}
This structure differs from the equivalent expression (1.71) for the
normalization constant $Z_T^\varepsilon$ only by the two insertions of the
matrix elements
\begin{equation}
  \bra{\tilde{\Phi}_0}O_1\ket{\Phi_0}
  \qquad \text{and} \qquad
  \bra{\tilde{\Phi}_{n_t}}O_2\ket{\Phi_{n_t}}.
\end{equation}

The operators $O_1$ and $O_2$ are expressions built from field operators and
their conjugate momenta. Two typical examples are
\begin{equation}
  O_A = \hat{\Phi}(n_0)^\dagger, \qquad
  O_B = \sum_{n\in\Lambda^3}\hat{\Phi}(n)\,e^{-ianp},
\end{equation}
where the first operator $O_A$ creates from the vacuum a field quantum at some
position $n_0$, and the operator $O_B$ annihilates a quantum projected to
momentum $p$. All operators $O_1$ and $O_2$ that can appear in (1.78) can be
written as expressions of field operators and conjugate momenta (1.55). Thus,
we obtain for matrix elements (1.79) of combinations of field operators
\begin{equation}
  \bra{\tilde{\Phi}}O\ket{\Phi} = O[\Phi]\,\delta(\tilde{\Phi} - \Phi),
\end{equation}
where we have made use of (1.57) and (1.58). The object $O[\Phi]$ is no longer
an operator but a functional of the classical field variables $\Phi$. A
functional maps a field configuration $\Phi$, specified by the set of all field
values $\Phi(n)$, $n \in \Lambda^3$, into the complex numbers. For the two
operators $O_A$ and $O_B$ in the example of (1.80) above, these functionals
read ($*$ denotes complex conjugation)
\begin{equation}
  O_A[\Phi] = \Phi(n_0)^*, \qquad
  O_B[\Phi] = \sum_{n\in\Lambda^3}\Phi(n)\,e^{-ianp}.
\end{equation}

Inserting the expression (1.81) for the matrix elements of the operators into
(1.78) and integrating out the fields $\tilde{\Phi}_0$ and $\tilde{\Phi}_{n_t}$
we find
\begin{equation}
  \ev{O_2(t)\,O_1(0)}_T^\varepsilon
  = \frac{C^{N^3N_T}}{Z_T^\varepsilon}\int D[\Phi]\;e^{-S_E[\Phi]}\,
    O_2[\Phi(\cdot,n_t)]\,O_1[\Phi(\cdot,0)].
\end{equation}
In the last step we have collected and combined the different terms, just as we
did for the partition function. In particular, we obtain the same expression
for the action $S_E$ as given in (1.74) and use again the measure $D[\Phi]$ as
defined in (1.75). The functionals $O_1[\Phi(\cdot,0)]$ and $O_2[\Phi(\cdot,n_t)]$
are the lattice transcriptions of the original operators $O_1$ and $O_2$ acting
in Hilbert space. The functional $O_1[\Phi(\cdot,0)]$ is evaluated for the
fields $\Phi(\cdot,0)$ with time argument 0, while the functional
$O_2[\Phi(\cdot,n_t)]$ is evaluated for the fields $\Phi(\cdot,n_t)$ with time
argument $n_t$, connected to the Euclidean time $t$ via $t = \varepsilon n_t$.
The spatial arguments of the fields, used for evaluating the functionals, were
replaced by a dot, since the functionals map the whole field configuration,
defined by the set of all field values at the given time, into the complex
numbers. We stress that the functionals $O_1[\Phi(\cdot,0)]$ and
$O_2[\Phi(\cdot,n_t)]$ as well as the Boltzmann factor $\exp(-S_E[\Phi])$ are
numbers, not operators, and therefore their relative order is no longer
crucial. Let us finally remark that the factor $C^{N^3N_T}$ in (1.83) is the
same factor that also shows up in the expression (1.76) for the partition
function. The two factors in the denominator and the numerator cancel and we
can omit them from now on.

\section{Quantization with the path integral}

In the last section we completed the derivation of the path integral
representation of the Euclidean correlators $\ev{O_2(t)\,O_1(0)}_T^\varepsilon$.
In this section we now discuss several aspects of the equations we have derived
and present some general concepts and strategies of lattice field theory. In
particular, we present the idea of directly using the path integral as a method
for quantizing a system. For further reading on this subject we recommend
[7, 10, 11].

\subsection{Different discretizations of the Euclidean action}

Let us begin our discussion of the path integral with a closer look at the
expression (1.74) for the Euclidean action. When inspecting the derivative
terms -- the first two terms on the right-hand side of (1.74) -- one notes a
difference between the two discretizations of the derivatives. The
discretization for the temporal derivative uses forward differences, whereas
for the spatial derivatives central differences are used. While the
discretization of the spatial derivative was chosen in (1.50), the temporal
derivative is a result of the stepwise transport in Euclidean time with the
Trotter formula. Thus it contains only nearest neighbors. Such a discretization
has, however, larger discretization errors than using central differences. From
the Taylor expansion $f(x \pm \varepsilon) = f(x) \pm \varepsilon f'(x)
+ \varepsilon^2 f''(x)/2 \pm \varepsilon^3 f'''(x)/6 + \cdots$ one finds that
for the forward differences
\begin{equation}
  \frac{f(x+\varepsilon) - f(x)}{\varepsilon} = f'(x) + \order{\varepsilon},
\end{equation}
while for the discretization with central differences
\begin{equation}
  \frac{f(x+\varepsilon) - f(x-\varepsilon)}{2\varepsilon}
  = f'(x) + \order{\varepsilon^2}.
\end{equation}
Due to the smaller discretization errors it is more advantageous to use the
central difference formula also for the temporal derivative.

A second change that is usually implemented is to use the same temporal and
spatial lattice constant, and to set
\begin{equation}
  \varepsilon = a.
\end{equation}
Although $\varepsilon$ and $a$ have a different origin, both of them should be
small: The spatial lattice constant determines down to which length scale we
can resolve the system, while the temporal constant determines how well the
Trotter formula approximates the Euclidean time transporter. When using the
central differences for time and implementing (1.86), the Euclidean action
reads
\begin{equation}
  S_E[\Phi] = a^4\sum_{n\in\Lambda}\Big[
    \frac{1}{2}\sum_{\mu=1}^4
      \Big(\frac{\Phi(n+\hat{\mu}) - \Phi(n-\hat{\mu})}{2a}\Big)^2
    + \frac{m^2}{2}\Phi(n)^2 + V(\Phi(n))\Big].
\end{equation}
It is obvious that in this form the action treats the spatial directions and
Euclidean time equally. Up to the fact that for Euclidean space we have no
relative sign between temporal and spatial components of the metric tensor, the
equal footing of space and time in (1.87) is the same as in the action (1.41)
and (1.42) with the Minkowski metric. The metric tensor which is used in (1.42)
to transform between covariant and contravariant indices is replaced by a
Kronecker delta.

We finally remark that in Minkowski space the components run from $\mu = 0$ to
3, where $\mu = 0$ refers to time. To stress that Euclidean metric is used the
index $\mu$ in (1.87) runs from 1 to 4, and $\mu = 4$ refers to Euclidean time.

\subsection{The path integral as a quantization prescription}

The action with central differences for time is not derived from first
principles in the way presented above. The changes we performed in the action
signals that we now change the philosophy for the quantization of the system.
Let us discuss this new quantization prescription.

In the derivation of the last section the quantization of the system entered by
enforcing the canonical commutation rules (1.52) and (1.53) for the fields and
their conjugate momenta. Subsequently, the expression
\begin{equation}
  \ev{O_2(t)\,O_1(0)}_T = \frac{1}{Z_T}\int D[\Phi]\;e^{-S_E[\Phi]}\,
    O_2[\Phi(\cdot,n_t)]\,O_1[\Phi(\cdot,0)],
\end{equation}
for the Euclidean correlators as path integrals was derived (from now on we
drop the superscripts $\varepsilon$ that indicate the discretization of time).
Also the partition function is given as a path integral
\begin{equation}
  Z_T = \int D[\Phi]\;e^{-S_E[\Phi]}.
\end{equation}
The measure
\begin{equation}
  D[\Phi] = \prod_{n\in\Lambda}d\Phi(n),
\end{equation}
is a product measure of integration measures for the classical field variables
at all points $n$ of the 4D lattice $\Lambda$.

Our new quantization prescription, which we introduce now, is based directly on
Eqs. (1.88), (1.89), and (1.90). It states that the quantization is no longer
implemented by enforcing canonical commutation rules for the commutators, but
instead by a path integral over classical field variables. The steps involved
in the quantization of a system are
\begin{itemize}
  \item \textbf{Step 1:} The continuous space-time is replaced by a 4D
        Euclidean lattice with lattice constant $a$. The degrees of freedom are
        classical field variables $\Phi$ living on the lattice.
  \item \textbf{Step 2:} The Euclidean action $S_E[\Phi]$ of the system is
        discretized on the lattice such that in the limit $a \to 0$ the
        Euclidean continuum action is obtained. The Boltzmann weight factor is
        $\exp(-S_E[\Phi])$.
  \item \textbf{Step 3:} The operators that appear in the Euclidean correlator
        one wants to study are translated to functionals by replacing the field
        operators with the classical lattice field variables.
  \item \textbf{Step 4:} Euclidean correlation functions are computed by
        evaluating these functionals on some lattice field configuration,
        weighting them with the Boltzmann factor and integrating over all
        possible field configurations.
\end{itemize}

In the next chapters we will implement this prescription for QCD. In this
approach there are, however, several different possibilities from which one can
choose. We have already seen that various discretizations of derivatives can be
used which differ in their discretization errors. An important question is
whether different discretizations of a system all lead to a theory where the
corresponding Euclidean correlators $\ev{O_2(t)\,O_1(0)}_T$ can be interpreted
according to our first key equation
\begin{equation}
  \lim_{T\to\infty}\ev{O_2(t)\,O_1(0)}_T
  = \sum_n \bbraket{0}{O_2}{n}\,\bbraket{n}{O_1}{0}\,e^{-tE_n},
\end{equation}
such that energy levels and physical matrix elements can be extracted. In other
words: Does a given lattice discretization of the Euclidean action used in a
path integral give rise to a proper quantum mechanical Hilbert space? If yes,
how can the $n$-point functions with Minkowski metric be reconstructed from the
Euclidean correlators? A very general answer to these questions is given by the
so-called Osterwalder--Schrader reconstruction (see, e.g., [12] for a textbook
presentation). If the Euclidean correlators calculated from the path integral
obey a certain set of axioms, then the Hilbert space for the Minkowski theory
can be obtained in a constructive way. For the Wilson formulation of lattice
QCD, which we will present in Chaps. 2 and 5, this construction was generalized
to lattice QCD in [13] (see also [14]). For an alternative construction on the
lattice, based on the transfer matrix, the interested reader is referred to
[15]. For our presentation here it is sufficient to conclude that we have some
freedom in the discretization of QCD, and when the axioms are obeyed the
physical Hilbert space can be constructed.

Once physical observables are calculated on the lattice one is interested in
their values in the limit $a \to 0$, the so-called continuum limit. This limit
can be understood in several ways. We have already encountered the so-called
naive or classical continuum limit when we required the lattice discretization
of the action to approach its continuum counterpart for $a \to 0$. However, the
naive continuum limit for the action is only used as a guiding principle in the
construction of the lattice theory. The fully quantized theory requires the
evaluation of the path integral, giving rise to results for observables which
are involved functions of $a$. In actual numerical calculations the continuum
limit is approached in a different way: The couplings of the theory are driven
to their critical values, i.e., to the values where the system undergoes a
phase transition. When doing so, physical scales (e.g., the proton size) become
large in lattice units. In other words, the resolution of our lattice becomes
finer and finer. Certainly these few remarks are too short to understand this
procedure which we will discuss in great detail in Sect. 3.5. This true
continuum limit is based on concepts from statistical mechanics and we conclude
this chapter by discussing the relation of lattice field theory to statistical
mechanics.

\subsection{The relation to statistical mechanics}

As announced, there is a structural equivalence between our field theory on the
lattice and statistical mechanics. A prototype system of statistical mechanics
is a spin system. The degrees of freedom are classical spin variables $s_n$
located on some lattice which is typically 3D. The energy of the system is a
functional $H[s]$ of the spins. In the canonical ensemble, i.e., for the system
in a heat bath with temperature $T$, the probability $P[s]$ of finding the
system in a particular configuration $s$ is given by
\begin{equation}
  P[s] = \frac{1}{Z}\,e^{-\beta H[s]},
\end{equation}
where $\beta$ is the inverse temperature $\beta = 1/k_B T$, with $k_B$ denoting
the Boltzmann constant. The partition function $Z$ is given by
\begin{equation}
  Z = \sum_{\{s\}}e^{-\beta H[s]},
\end{equation}
where the sum runs over all possible spin configurations $\{s\}$. The
expectation value of some observable $O$ is given by
\begin{equation}
  \ev{O} = \frac{1}{Z}\sum_{\{s\}}e^{-\beta H[s]}\,O[s].
\end{equation}

The similarity between (1.88) and (1.94) is obvious. The Boltzmann factor
$\exp(-\beta H)$ is replaced by the weight factor $\exp(-S_E)$ and the
summation over all spin configurations by the integral over configurations of
the classical field $\Phi$. Thus, the name ``partition function'' for the
normalization factor of the path integral now appears natural. The structural
equivalence between lattice field theory and statistical mechanics is an
important observation also from a technical point of view: It allows one to
apply analytical and numerical methods, originally developed in statistical
mechanics, to lattice field theory. Some of these techniques will be discussed
in subsequent chapters.

\section*{References}

\begin{enumerate}[label={[\arabic*]}]
  \item C.~Cohen-Tannoudji, B.~Diu, and F.~Lalo\"e: \emph{Quantum Mechanics}
        (John Wiley and Sons, New York, London, Sidney, Toronto 1977).
  \item J.~J.~Sakurai: \emph{Modern Quantum Physics} (Addison-Wesley, Reading,
        Massachusetts 1994).
  \item A.~Messiah: \emph{Quantum Mechanics}, Vols.~1 and 2 (Dover
        Publications, Mineola, New York 1958).
  \item J.~W.~Negele and H.~Orland: \emph{Quantum Many Particle Systems}
        (Westview Press -- Advanced Book Classics, Boulder 1998).
  \item E.~Prugove\v{c}ki: \emph{Quantum Mechanics in Hilbert Space}, 2nd ed.
        (Academic Press, New York, London, Toronto, Sydney, San Francisco
        1981).
  \item M.~Reed and B.~Simon: \emph{Methods of Modern Mathematical Physics},
        revised and enlarged ed., Vol.~1: Functional Analysis (Academic Press,
        Boston, San Diego, New York, London, Sydney, Tokyo, Toronto 1980).
  \item G.~Roepstorff: \emph{Path Integral Approach to Quantum Physics}
        (Springer, Berlin, Heidelberg, New York 1996).
  \item C.~Itzykson and J.-B.~Zuber: \emph{Quantum Field Theory} (McGraw-Hill,
        New York 1985).
  \item M.~E.~Peskin and D.~V.~Schroeder: \emph{An Introduction to Quantum
        Field Theory} (Addison-Wesley, Reading, Massachusetts 1995).
  \item R.~P.~Feynman and A.~Hibbs: \emph{Quantum Mechanics and Path
        Integrals} (McGraw-Hill, New York 1965).
  \item R.~J.~Rivers: \emph{Path Integral Methods in Quantum Field Theory},
        Cambridge Monographs on Mathematical Physics (Cambridge University
        Press, Cambridge, New York 1990).
  \item J.~Glimm and A.~Jaffe: \emph{Quantum Physics. A Functional Integral
        Point of View}, 2nd ed. (Springer, New York 1987).
  \item K.~Osterwalder and E.~Seiler: Ann.\ Phys.\ 110, 440 (1978).
  \item E.~Seiler: \emph{Gauge Theories as a Problem of Constructive Quantum
        Field Theory and Statistical Mechanics}, Lect.\ Notes Phys.\ 159.
        Springer, Berlin, Heidelberg, New York (1982).
  \item M.~L\"uscher: Commun.\ Math.\ Phys.\ 54, 283 (1977).
\end{enumerate}
